\abresubcapitulo
  {Tempo, tecnologia e trabalho}
  {Redes técnicas, produtividade e disputa pelo tempo na escala $6\times1$}
  {GEOGRAFIA \;\textbullet\; MATERIAL COMPLEMENTAR \;\textbullet\; ENSINO MÉDIO / ENEM}

\begin{questaocentral}
\textbf{QUESTÃO CENTRAL:} Se as tecnologias permitem produzir mais em menos tempo, por que esse aumento de produtividade não se transforma automaticamente em redução da jornada e ampliação do tempo livre?
\end{questaocentral}

\subsection*{Objetivos de aprendizagem}

Ao final deste subcapítulo, você deverá ser capaz de:

\begin{itemize}
  \item interpretar uma paisagem urbana como expressão de desigualdades sociais e econômicas;
  \item compreender, de maneira introdutória, o Meio Técnico-Científico-Informacional;
  \item reconhecer a distribuição desigual do comando, dos ganhos e dos riscos nas redes produtivas;
  \item explicar por que a redução da jornada não decorre automaticamente do desenvolvimento tecnológico.
\end{itemize}

\section{Da paisagem à rede geográfica}

\begin{exemploconcreto}
\textbf{SITUAÇÃO INICIAL -- AVENIDA PAULISTA:} Imagine um trabalhador sob chuva, cercado por sedes empresariais, bancos, redes de comunicação e intensa circulação de mercadorias. Ele afirma que trabalha muito, mas não consegue transformar seu esforço em segurança ou expectativa estável de futuro. O contraste não é apenas social: ele aparece materializado na própria paisagem urbana.
\end{exemploconcreto}

A paisagem reúne, no mesmo lugar, posições muito diferentes dentro da economia. A proximidade espacial entre riqueza e precariedade não significa que todos tenham o mesmo poder sobre as redes que atravessam a cidade. Para compreender essa desigualdade, precisamos olhar além dos aparelhos e observar o conjunto de infraestruturas, empresas, informações e trabalhadores que permite à atividade econômica funcionar.

\begin{conceito}
\textbf{CONCEITO-CHAVE -- MEIO TÉCNICO-CIENTÍFICO-INFORMACIONAL:} forma de organização do espaço geográfico marcada pela integração entre técnica, ciência, informação, infraestruturas, empresas e ações humanas. Um aplicativo, por exemplo, só funciona porque está conectado a redes de comunicação, pagamentos, GPS, vias urbanas, estabelecimentos e trabalhadores \citep{santos}.
\end{conceito}

\begin{figure}[htbp]
\centering
\begin{tikzpicture}[
  node distance=11mm and 15mm,
  box/.style={rounded corners=2pt,draw=petroleo,thick,minimum width=3cm,minimum height=9mm,align=center,font=\small\bfseries},
  flow/.style={<->,>=Stealth,semithick,draw=cinza},
  label/.style={font=\scriptsize,color=cinza,align=center}
]
  \node[box,fill=azulclaro] (consumidor) {Consumidor};
  \node[box,fill=petroleo,text=white,right=of consumidor] (plataforma) {Plataforma};
  \node[box,fill=ouroclaro,above right=8mm and 28mm of plataforma] (estabelecimento) {Estabelecimento};
  \node[box,fill=verdeclaro,below right=8mm and 28mm of plataforma] (trabalhador) {Trabalhador};

  \draw[flow] (consumidor) -- node[label,above]{pedido · pagamento · dados} (plataforma);
  \draw[flow] (plataforma) -- node[label,above,sloped]{demanda · informação} (estabelecimento);
  \draw[flow] (plataforma) -- node[label,below,sloped]{preço · avaliação} (trabalhador);

  \node[draw=cinza!35,fill=cinzaclaro,rounded corners=2pt,font=\scriptsize,text=petroleoescuro,
        below=18mm of plataforma,minimum width=12.4cm,minimum height=8mm]
        {Infraestruturas: internet · GPS · pagamentos · vias · transporte};
\end{tikzpicture}
\caption{Rede simplificada do trabalho organizado por plataforma.}
\fontefigura{elaboração didática para este protótipo.}
\end{figure}

\begin{frasechave}
\textbf{FRASE-CHAVE:} A rede funciona digitalmente, mas depende de corpos trabalhando em lugares concretos.
\end{frasechave}

\begin{atividade}
\textbf{LEITURA DO ESQUEMA:} Quais elementos materiais são indispensáveis para que essa rede aparentemente ``digital'' funcione? Quem pode definir preços e critérios? Quem permanece mais exposto aos riscos do trânsito e da espera?
\end{atividade}

\section{Tecnologia também reorganiza o tempo}

As inovações técnicas ampliam possibilidades, mas não determinam previamente como elas serão utilizadas. A iluminação artificial permitiu que fábricas, lojas e serviços prolongassem suas atividades para além do tempo natural do dia. Em vez de reduzir automaticamente a jornada, uma tecnologia pode ampliar o período disponível para produzir.

\begin{conceito}
\textbf{PRODUTIVIDADE:} relação entre os resultados obtidos e os recursos ou o tempo empregados. Um aumento de produtividade permite produzir mais no mesmo intervalo, mas não define quem receberá os ganhos nem quanto tempo será destinado ao descanso.
\end{conceito}

\subsection*{Questão de fixação}

\begin{questao}
\textbf{QUESTÃO AUTORAL:} A difusão da iluminação artificial ampliou o funcionamento noturno de fábricas e estabelecimentos. Do ponto de vista das relações entre técnica e trabalho, esse processo demonstra que:

\begin{enumerate}[label=\Alph*),leftmargin=2em]
  \item a inovação tecnológica elimina conflitos em torno da duração da jornada.
  \item o desenvolvimento técnico conduz necessariamente ao aumento do tempo livre.
  \item a técnica possui efeitos idênticos para todos os grupos sociais.
  \item o uso social de uma tecnologia pode ampliar o tempo disponível para a produção.
  \item a organização do trabalho independe das relações políticas e econômicas.
\end{enumerate}
\end{questao}

\begin{resolucao}
\textbf{RESOLUÇÃO COMENTADA -- D:} a iluminação artificial cria a possibilidade de estender a atividade produtiva. A duração da jornada, porém, depende de relações sociais, direitos e decisões políticas; não é consequência automática da tecnologia.
\end{resolucao}

\begin{notaeditorial}
\textbf{NOTA DO PROTÓTIPO:} o plano original menciona ``ENEM 2015, questão 02'' na sequência didática e ``Questão 03'' nas referências. A equipe deverá conferir o caderno antes de reproduzir a questão oficial. Esta caixa existe para demonstrar o tratamento de uma pendência editorial e será removida da edição destinada aos estudantes.
\end{notaeditorial}

\section{Divisão territorial do trabalho}

\begin{exemploconcreto}
\textbf{ESTUDO DE CASO -- RANA PLAZA:} Em 2013, o desabamento do edifício Rana Plaza, em Bangladesh, matou 1.134 pessoas. O prédio abrigava fábricas de roupas integradas a cadeias internacionais. A tragédia tornou visível uma geografia na qual marcas e mercados consumidores podiam estar distantes dos locais onde o trabalho e o risco se concentravam \citep{oxfam}.
\end{exemploconcreto}

\begin{conceito}
\textbf{DIVISÃO TERRITORIAL DO TRABALHO:} distribuição das atividades e funções econômicas entre diferentes lugares e posições de uma mesma rede. Alguns espaços concentram comando, tecnologia, informação e lucros; outros concentram execução, remunerações menores e exposição mais direta aos riscos.
\end{conceito}

\begin{frasechave}
\textbf{FRASE-CHAVE:} A mesma rede que conecta os lugares também distribui de forma desigual comando, ganhos, trabalho e risco.
\end{frasechave}

\section{Cadeias globais e plataformas digitais}

Uma cadeia global de roupas e o trabalho organizado por plataformas não são fenômenos equivalentes. A comparação serve para localizar uma questão comum: quem controla aspectos centrais da rede e quem assume diretamente seus custos e riscos?

\begin{table}[htbp]
\centering
\small
\renewcommand{\arraystretch}{1.22}
\begin{tabularx}{\textwidth}{>{\bfseries}Y Y Y}
\toprule
\rowcolor{petroleo}
\color{white}Pergunta & \color{white}\bfseries Cadeia global de roupas & \color{white}\bfseries Trabalho por plataformas \\
\midrule
Quem controla aspectos centrais? & Marcas, contratantes e gestores da cadeia & Empresa proprietária da plataforma \\
\rowcolor{cinzaclaro}
Quem executa o trabalho? & Trabalhadores das fábricas & Entregadores e motoristas \\
Quem assume custos e riscos diretos? & Trabalhadores e fornecedores locais & Trabalhador \\
\bottomrule
\end{tabularx}
\caption{Comparação analítica entre duas formas distintas de organização do trabalho.}
\end{table}

\begin{conceito}
\textbf{TRANSFERÊNCIA DE RISCOS:} processo pelo qual uma empresa mantém controle sobre aspectos centrais da atividade, mas desloca parte dos custos e riscos para trabalhadores ou fornecedores. No trabalho por plataformas, isso pode incluir veículo, combustível, manutenção, tempo de espera e risco de acidente.
\end{conceito}

\begin{atividade}
\textbf{COMPARE SEM IGUALAR:} identifique uma semelhança e uma diferença entre a cadeia global de roupas e o trabalho por plataformas. Sua resposta deve mencionar comando, execução ou risco.
\end{atividade}

\section{Produtividade e escala \texorpdfstring{$6\times1$}{6x1}}

Tecnologias podem acelerar a circulação, conectar lugares e ampliar a capacidade de produzir. Nenhuma dessas possibilidades decide, por si mesma, quem receberá os ganhos de produtividade, quantos dias serão trabalhados ou quanto tempo será destinado ao descanso.

\begin{frasechave}
\textbf{A escala $6\times1$ não é uma necessidade criada pela tecnologia.} Ela é uma forma social e política de organizar e distribuir o tempo de trabalho.
\end{frasechave}

\section*{Síntese}

\begin{questaocentral}
\textbf{RETOMANDO A QUESTÃO CENTRAL:} a tecnologia modifica a produção e reorganiza o espaço e o tempo, mas a transformação dos ganhos de produtividade em segurança, descanso ou redução da jornada depende de direitos, decisões políticas, organização coletiva e relações de poder.
\end{questaocentral}

\begin{atividade}
\textbf{EM UMA FRASE:} complete: ``A tecnologia modifica a produção, mas o tempo livre depende de...'' Depois, justifique sua resposta com um exemplo do subcapítulo.
\end{atividade}

\subsection*{Vocabulário de revisão}

\begin{itemize}[label={},leftmargin=0pt,itemsep=4pt]
  \item \textbf{\color{petroleoescuro}Meio Técnico-Científico-Informacional:} integração entre técnica, ciência, informação, infraestruturas e ações na organização do espaço.
  \item \textbf{\color{petroleoescuro}Divisão territorial do trabalho:} distribuição de funções econômicas entre diferentes lugares e posições de uma rede.
  \item \textbf{\color{petroleoescuro}Transferência de riscos:} deslocamento de custos e riscos para participantes com menor controle sobre a atividade.
  \item \textbf{\color{petroleoescuro}Produtividade:} relação entre resultados produzidos e recursos ou tempo empregados.
\end{itemize}
